% 本文件由 LaTeX 排版 Agent 根据有效 translation.json 自动生成。
% author=Miscellaneous; work_id=0515; title=论重圣洗

\chapter{论重圣洗}

% structure: p0001 source=h1 target=omitted reason=page-title-already-rendered-by-parent

\Emph{论点：凡曾因主耶稣基督之名受圣洗者，不应再受圣洗。}\par

1. \Emph{我注意到}，在弟兄们中间曾有人问及，对于某些人——他们虽在异端中受圣洗，却仍因我们的主耶稣基督之名受圣洗，后来离开异端，作为恳求者投奔天主的教会，全心悔改，直到现在才认清自己错误的定谳，并向教会恳求救恩的帮助——应当特别采取何种做法。\Emph{问题是}：按照最古老的习俗和教会传统，对于那些在\Emph{教会}之外，却仍因我们的主耶稣基督之名所接受的圣洗，是否只需由主教为他们覆手以领受圣神，而这一覆手礼就能为他们提供更新而完美的信德印记就足够了？或者，他们是否必须重复圣洗，仿佛若不重新领受圣洗就一无所获，如同从未因耶稣基督之名受圣洗一样？因此，关于这个新问题，有人谈论了一些已写就的答复和回应，双方都以极大的热忱试图摧毁对方所写的内容。在我看来，在这种争论中，如果我们每个人都满足于所有教会的可敬权威，并以应有的谦卑不愿创新，看到没有任何反驳的余地，那么本不会产生任何争议或讨论。因为一切既可疑又模棱两可、且在\Emph{那些}明智而忠实之人的不同意见中被确立的事情，如果被认为违背了所有那些值得称颂的圣洁而忠实之人的古老、令人难忘且至为庄严的惯例，就理应受到谴责；因为在一件已经安排和规定的事情上，任何被提出而反对教会安宁与和平的东西，都只会导致不和、纷争与分裂。其中找不到任何别的果实，唯独这一点：有一个人，无论他是谁，会在某些善变的人中被虛荣地宣告为极具智慧和恒心；并且，他因具有异端者的傲慢——异端者毁灭中唯一的安慰是不显得独自犯罪——而在那些与他最相似、最相投的人中扬名，被认为纠正了所有教会的错误和缺点。因为所有异端者的愿望和目的就是尽可能多地编造这种诽谤，来反对我们至圣的母亲教会，并认为发现任何可以归咎于她的罪行甚至愚蠢之处，就是莫大的荣耀。既然没有一个心智健全的忠实信徒胆敢持有这样的观点，尤其是没有任何被授予任何圣职的人——更不用说主教职的人——敢这样做，那么主教们自己设计这样的丑闻，并且不顾法律和全部圣经的训诫，以他们自己的耻辱和危险，毫不敬畏地揭露他们的母亲教会的耻辱——如果他们以为这其中有任何耻辱的话——这简直像是一件怪事。尽管在这件事上，教会除了这些人本身的错误外，并无耻辱。因此，对于这种人来说，更为严重的罪过是：如果他们在最古老的惯例中所指责的（仿佛做得不对）显然且有力地表明，既被我们先人正确遵守，也被我们正确遵守；因此，即使我们以双方均等的论据进行辩论，但由于创新之事若不引起弟兄们之间的不和及对教会的损害就无法确立，那么，无论如何——不论对错，正如他们所说，即违背善与正当——都不应轻率地将其像污点一样抛向我们母亲教会；这种鲁莽和不敬的耻辱理应归于那些尝试这样做的人。然而，按照宗徒的训诫，我们无权\BibleQuote{说同样的话，免得我们中间有分裂；} \BibleRef{格前 1:10} 然而，就我们所能，我们努力证明这一论点的真实状况，并劝告那些好争辩的人，即使是现在，也要管好自己的事，因为如果他们最终听从这良善的劝告，我们将会取得很大的成果。因此，我们将根据需要，将与此主题相关的所有圣经经文汇集为一整篇。我们将尽可能清楚地协调那些看似不同或含义各异的部分；我们将以我们微薄的能力探讨每种方法的效用和益处，以便向所有弟兄推荐：在教会中采取最有益安康的形式与和平的习俗。\par

2. 因此，对于那些前来讨论拯救性且现代的，即属于圣神与福音的洗礼的人，首先出现的是一位普遍极为熟悉的宣告，由\Person{若翰洗者}所宣告并开始。他略微离开了律法，即最古老的\Person{梅瑟}洗礼，为新的真恩宠预备道路，既以他当时所施行的水和悔改的洗礼逐渐占据犹太人的耳朵，又以对将来属神洗礼的宣告占据他们的心，劝诫他们说：\Quote{那在我以后来的，比我更有能力，我连给祂解鞋带也不配；祂要以圣神和火给你们施洗。}\BibleRef{玛 3:11} 正因如此，我们也应从此处开始这篇论述。因为在\WorkTitle{宗徒大事录}中，主复活后，证实了若翰的这话，\Quote{吩咐他们不要离开\Place{耶路撒冷}，但要等候圣父所应许的，就是你们从我听说的，\Emph{祂说}；因为若翰确实用水施了洗，但不多几日以后，你们要因圣神受洗。}\BibleRef{宗 1:4} \Person{伯多禄}在向宗徒们为自己辩护时，也叙述了主的这些话说：\Quote{我一开口讲，圣神就降在他们身上，像当初降在我们身上一样；我就想起主的话说：若翰用水施洗，但你们要因圣神受洗。所以，如果他们得了与我们相同的恩赐，我们因信主耶稣基督，我是谁，竟能阻止天主吗？}\BibleRef{宗 11:15} 又说：\Quote{诸位同人及弟兄们，你们知道，昔日天主在我中间选择了，叫外邦人从我口中听到福音的话而相信。认识人心的天主，给他们作了证，赐给他们圣神，如同赐给我们一样。}\BibleRef{宗 15:7} 因此，我们应当思考这话的力量与权能。因为主对那些后来因信而将要受洗的人说，他们必须受洗，不是像祂用水那样为悔改，而是因圣神受洗。对于这宣告，我们中任何人都不能怀疑，那么人们是因圣神受洗的原则就显而易见了。因为信徒尤其只在圣神本身内受洗。因为若翰分别开来，说他确实用水施洗，但有一位要来，祂要因圣神施洗，借着天主的恩宠和大能；他们借着\Emph{圣神的}赐予和隐秘效果的运作而如此。此外，他们在圣神和水的洗礼中也同样如此。此外，每一个人在自己血中的洗礼中也是如此。正如圣经向我们宣告的那样，我们将从圣经中为我们将要叙述的每一具体事例提供明确的证据。\par

3. 对此，你或许，那位带来某些新奇事物的人，可能立即不耐烦地回答说，正如你所习惯的那样，主在福音中说：\Quote{人若不重生于水和圣神，就不能进天主的国。}由此可见，唯有在其中圣神能居住的洗礼才是有益的；因为主自己受洗时，圣神曾降临在祂身上，而且祂的言行完全一致，这样的奥迹不能与其他原则共存。对此，我们中没有人会如此愚昧或固执，敢于违背正义或真理提出反对，例如，对我们完整地、并在教会中以各种方式行事，以及按照纪律秩序永久遵守这些事的做法提出反对。但如果在同一新约中，我们在那件事上所遇到的相关事物有时被发现以某种方式分开、分离、排列、安排，就好像它们是单独的一样；让我们看看这些单独的实例本身有时是否不是不完整的，而是似乎完整且完备的。因为当通过主教的按手，圣神赐给每个相信的人时，如在\Place{撒玛黎雅}人的情形中，\Person{斐理伯}施洗之后，宗徒们通过按手在他们身上，就这样也把圣神赐给了他们。而为了使这事成为可能，他们自己为他们祈祷，因为那时圣神还没有降临在他们任何人身上，他们只是已经因主耶稣的名受了洗。此外，我们的主在复活后，向宗徒们吹了气，并对他们说：\Quote{你们领受圣神吧}\BibleRef{若 20:22}，唯有这样和仅仅这样把圣神赐给了他们。\par

4. 若情况如此，我的兄弟，你认为如何？如果一个人并非由主教施洗，以至于未能立即领受覆手，而在领受圣神之前去世，你判断他是否获得了救恩？因为，事实上，宗徒们和门徒们自己，他们也为他人施洗，并由主施洗，却没有立刻领受圣神，因为那时圣神尚未赐下，因为耶稣尚未受光荣。祂复活后，过了相当长的时间，圣神才降临——正如撒玛黎雅人一样，当他们由斐理伯施洗时，\Emph{并未领受恩赐}，直到宗徒们从耶路撒冷被邀请到撒玛黎雅，下去为他们覆手，藉覆手将圣神赐予他们。因为在那段时间里，他们中任何一位尚未获得圣神的人，都可能被死亡截断，至死被剥夺圣神的恩宠。同样，无疑今日这类事也常见，且频繁发生：许多人领洗后，未得主教覆手便离世，却仍被视为完美的信徒。正如埃塞俄比亚的太监，从耶路撒冷返回，诵读依撒意亚先知书，心存疑惑，在圣神的感动下从执事斐理伯那里听到真理，相信并受了洗；他从水中上来时，主的神将斐理伯提去，太监再也看不见他。因为他欢欢喜喜地继续前行，尽管如你所见，主教没有为他覆手使他领受圣神。但若你承认这点，并相信这是得救的，且不反驳所有信徒的意见，那么你必须承认：正如这一原则得以更广泛地讨论，另一原则也可以更广泛地确立——即，仅凭主教的覆手，因先有因我们的主耶稣基督之名所行的圣洗，圣神也可以赐予另一个悔改并相信的人。因为圣经已肯定，凡信基督的人，必须受圣神的洗，好使这些人也不被视为比那些完美的基督徒缺少什么；免得有必要追问，他们因耶稣基督之名所得的圣洗究竟是何种圣洗。除非，或许在先前的讨论中，关于那些只以耶稣基督之名受洗的人，你断定他们即使没有圣神也能得救，或者圣神不惯于只以这种方式赐予，而必须通过主教的覆手；或者甚至说，并非只有主教能赐予圣神。\par

5. 若果真如此，且这些事中任何一件都不能剥夺相信者的救恩，那么你自己也承认：信仰奥迹在某种程度上被分割，而且如你所争辩的，在必需介入时未能完成，这并不能剥夺一个相信并悔改之人的救恩。或者若你说这样的人不能得救，那我们便剥夺了所有主教的救恩，因为你以此种几乎确定的危险，使他们自己必须为所有生活在他们管辖下、分散各地、身体虚弱的人提供帮助，因为其他等级较低的圣职人员不敢施行同样的恩惠；这样，那些未能领受此恩惠而去世之人的血，就必须向主教追讨。此外，如你所知，我们发现，主有时无需水洗便将圣神赐予相信的人，正如\BibleBook{宗徒}{大事录}所记载：\BibleQuote{伯多禄还在讲这些话的时候，圣神降在所有听道的人身上。那些同伯多禄一起来的受割损的信徒，都惊讶圣神的恩赐也倾注在外邦人身上，因为他们听见他们说各种语言，并光荣天主。那时，伯多禄回答说：这些人既与我们一样领受了圣神，谁能阻止他们受水洗呢？就吩咐他们因耶稣基督的名受洗。} \BibleRef{宗 10:44} 正如伯多禄后来更充分地教导我们关于同一外邦人的事，他说：\BibleQuote{藉着信德洁净了他们的心，并没有区别我们和他们。} \BibleRef{宗 15:9} 毫无疑问，人可以不藉水而受圣神的洗——正如你注意到，这些人在受水洗之前已经领受了圣神；这样，若翰和我们主自己的宣告都得到了满足——因为他们既无宗徒的覆手，也无洗礼池，却获得了应许的恩宠，他们后来才领受了这些。他们的心被洁净，天主同时因他们的信德赐予他们罪赦；因此后来的圣洗只给了他们这个好处：他们领受了耶稣基督之名的呼求，使他们的侍奉和信德的完整无所欠缺。\par

6. 此外——从本讨论的另一面来看——我们主的那些门徒，他们先前已受过洗，在五旬节日圣神终于降临到他们身上，确实是由天主旨意从天上降下，并非出于自发，而是为这职务而倾注，并且降临到他们每人身上。他们虽然已是义人，且如我们所说，已受过主的洗礼，甚至宗徒们也是如此，然而在祂被捕的那夜，他们都离弃了祂。甚至伯多禄本人，他曾夸口要持守信德，最顽固地抗拒主自己的预言，但最终否认了祂，借此向我们表明：他们在其间以任何方式所犯的一切罪，后来因着真诚的信德，无疑都藉圣神的洗礼被除去了。并且，我认为，宗徒们之所以吩咐他们在圣神内所劝告的人，要因基督耶稣之名受洗，别无其他原因，除非是藉洗礼呼号耶稣之名于某人身上，能为他获得救恩带来不少益处，正如伯多禄在\WorkTitle{宗徒大事录}中所说：\BibleQuote{因为在天下人间，没有赐下别的名字，使我们赖以得救。} \BibleRef{宗 4:12} 同样，宗徒保禄也阐明，显示天主高举了我们的主耶稣，并\BibleQuote{赐给祂一个名字，超越一切名字，为使因耶稣之名，天上、地上和地下的一切，都屈膝叩拜，并一切唇舌都承认耶稣基督是主，以光荣天主圣父。} 那人受洗时，若因耶稣之名呼求于他，尽管他可能在某种错误下领受洗礼，仍不会阻碍他日后某个时候认识真理、改正错误、来到教会和主教面前，并在人前真诚宣认我们的耶稣；这样，当主教给他覆手时，他也能领受圣神，并且不会失去那先前因耶稣之名的呼求。我们中无人可否认可此，尽管这呼求若孤立而自足，不足以为得救，免得我们据此原理相信，甚至异民和异端者中妄用耶稣之名者，无需真实完整的事物就能得救。然而，极有益处的是相信：这种呼号耶稣之名，连同改正错误、承认真理的信仰、并除去以往生活中的一切污点，若正当履行于这类人中，与天主的奥迹一起，便获得原本没有的地位；最终，在真信仰中，为保持标记的完整，并不构成阻碍，当所缺的补充加上时；并且这符合良好理由，符合如此多年、如此多教会、宗徒和主教的权威；正如对我们的至圣慈母教会第一次突然无理由地反叛此前如此多世纪以来的决定，是极大的不利和损害。因为并非其他原因，伯多禄——他早已受洗，且曾被主问及对主的看法，并因得自天父启示的真理，宣认基督不仅是我们主，而且是永生天主之子——后来被显示曾敌对那宣告自己受难的同一基督，因此被称为撒殚。\Emph{别无其他原因}，只因为将发生：有些人虽在自己判断上摇摆，在信仰和道理上有些跛行，尽管他们因耶稣之名受了洗，但若能在某段时间内纠正错误，并不因此被排除于救恩之外；反而在任何时候他们醒悟过来，藉悔改获得救恩的可靠希望，尤其当他们领受圣神——每个人当受祂的洗礼——时，他们就会如此意图。正如我们在\WorkTitle{福音}中看到，伯多禄不仅承受了这事，所有门徒都如此，主对他们——虽已受洗——后来却说：\BibleQuote{你们众人都要因我而跌倒。} \BibleRef{谷 14:27} 我们观察到，他们全都改正了信仰，在主复活后领受了圣神的洗礼。因此，我们今日也有理由相信：人们从先前的错误中改正后，可以领受圣神的洗礼，他们虽曾因主的名受水洗，但信仰可能有所欠缺。因为，一个人是否从未因我们的主耶稣基督之名受洗，或者他受水洗时有所跛行——这较不重要——只要后来在圣神的洗礼中显明对真理的真诚信德（这无疑更为重要），这是非常重要的。\par

7. 你也不可认为主所说的与这种处理相矛盾：\Quote{你们要去教训万民，因父、及子、及圣神之名给他们授洗。}\BibleRef{玛 28:19}因为，虽然这是真实正确的，并应在教会中以各种方式遵守，而且历来如此遵守，但我们仍需考虑，因着那名的尊荣与权能，耶稣之名的呼求不应被我们视为无效；以那名字，各样的权能惯常被执行，甚至有时被教会外的人所行。但基督的话有何意义呢？祂说，在审判之日，那些向祂说\Quote{主啊，主啊，我们不是因祢的名字预言过，因祢的名字驱逐过魔鬼，因祢的名字行过许多异能}的人，祂要否认、不认识他们，当时祂甚至加重语气回答他们说：\BibleQuote{我从来不认识你们，你们这些行不义的人，离开我去吧！}\BibleRef{谷 14:27}除非是要向我们表明，即使那些行不义的人，也能因基督之名的剩余能力而做这些善工？因此，耶稣之名的呼求应当被视为我们与所有其他人共有的主之奥秘的某种开端，这开端日后可由其余事物来完成。否则，这样的呼求若单独存留，便无功效，因为处于这种地位的人死后，不能给他增添任何事物，也不能补充，且在审判之日不能在任何事上帮助他；那时，他们将被主以我们上面提到的事责备，而这些人中，没有一个人在现世被任何人如此严酷残忍地禁止以我们上面指出的方式帮助自己。\par

8. 但你会像往常一样，反对我们，提出异议说：那些门徒受洗时，是完美而正确地受洗，不像这些异端者；而你必须基于他们的状况以及为他们施洗者的状况来作此假设。因此，我们回应你的这个主张，并非作为主门徒的控告者，而是出于不得已，因为我们必须通过理性探究救恩是在何时、何地、以何种尺度赐予我们每个人的。因为我们的主诞生了，祂就是基督，这一点门徒们有诸多理由相信，并非不公正：祂出生于犹大支派，达味家族，在白冷城；祂诞生时，天使同时向牧羊人宣告救主为他们诞生；祂的星在东方出现，被贤士们殷勤寻求、朝拜，并以辉煌的礼物和卓越的奉献尊荣祂；祂年幼时，坐在圣殿的法学士中间，智慧地辩论，令众人赞叹；祂受洗时，因天开、圣神降临并住在祂身上而受光荣，这是别人未曾有过的；此外，还有祂父亲和若翰洗者的见证；祂超越凡人的低下能力，洞察所有人的心与思念；祂以极大的权能医治弱症、恶习与疾病；祂以明证赦免罪过；祂一句话就驱逐魔鬼；祂一句话洁净癞病人；祂将水变酒，以奇妙的喜乐增广婚宴的欢庆；祂恢复或赐予盲人视力；祂以十足的自信维护父的教导；祂在旷野以五个饼使五千人饱足，剩下的碎块装满十二筐以上；祂按照仁慈随处复活死人；祂命令风与海平静；祂以脚行走在海面上；祂完全施行了所有奇迹。\par

9. 藉这些事，以及许多类似的归于祂光荣的事迹，似乎自然得出一个结论：无论犹太人如何思想基督，尽管他们不信我们的主耶稣基督，甚至连他们自己也认为这样一位如此伟大的人物，将不经死亡而永远存续，并永远拥有以色列国和整个世界的王权，且这王国不会被毁灭。因此，犹太人甚至胆敢用武力抓住祂，要膏立祂为王，而祂不得不逃避此事；于是，祂的门徒认为，除非祂首先亲自经历这短暂的生命延续到永恒的生命，否则祂不会以其他方式赐予他们永生。最后，当他们经过加里肋亚时，耶稣对他们说：\Quote{人子要被交在人手中，他们要杀害祂；三天以后祂要复活。}\BibleRef{谷 9:30}他们极其忧伤，因为正如我们所说，他们先前在心中和意念中形成了截然不同的想法。再者，这也是犹太人对祂反驳的话，当祂教导他们关于自己的事，并向他们预告将来的事时，他们说：\Quote{我们从法律上听说基督永远存留；你怎么说人子必须被举起来呢？}\BibleRef{若 12:34}因此，门徒心中对基督也有同样的假定，就连宗徒之长和领袖伯多禄本人也迸发出他不信的言论。因为当他与其他门徒一起被主问及他如何看待祂，即他认为祂是谁时，他首先承认了真理，说祂是基督、永生天主之子，并因此被祂称为有福，因为他不是按血肉，而是藉天父的启示获得了这真理；然而，同样是这位伯多禄，当耶稣开始向门徒指示祂必须往耶路撒冷去，受长老、司祭长和经师许多苦，并被杀害，第三天复活时，这位基督的真正宣认者却在数日后拉过祂，开始责斥祂说：\Quote{求祢怜悯自己，这事不会临到祢身上！}\BibleRef{玛 16:22}为此他理应听到主说：\Quote{撒殚，退到我后面去！你是我的绊脚石，因为你所体会的，不是天主的事，而是人的事。}这个对伯多禄的斥责在主被捕时变得更加明显，他因使女的惊吓而说：\Quote{我不知道你说什么，也不认识你。}\BibleRef{玛 26:70}并且又发着誓说了同样的话；第三次，他诅咒并发誓说不认识那个人，且不止一次，而是多次否认祂。这种心态，因为要持续到他直到主的受难，早已被主显明，好叫我们也不对此无知。再者，主复活后，祂的一个门徒克罗帕，在向主本人（如同对某个不认识的人）伤心地讲述所发生的事时，如同他所有同伴的错误一样，这样说道：\Quote{……关于纳匝肋人耶稣，祂是一位先知，在天主和全体百姓面前，行事说话都有权能。我们的司祭长及首领竟把祂交出去，判了死罪，钉在十字架上。但我们原希望祂就是那要拯救以色列的。}\BibleRef{路 24:20}除此之外，所有门徒也认为那些在复活后看见主的妇女们的报告是胡言乱语；他们中的一些人在亲眼看见祂后仍不相信，反而怀疑；而那些当时不在场的人则根本不信，直到后来被主亲自以各种方式责斥和责备；因为祂的死如此冒犯他们，以至于他们认为祂没有复活，而他们曾相信祂不应死去，因为与他们所相信的相反，祂确实死了。因此，就门徒本身而言，他们在我们提到的这些事上，被发现既没有纯正的信德，也没有完善的信德；而更为严重的是，他们竟然还给别人施洗，正如\BibleBook{若望}{福音}所记载的。\par

10. 此外，对于那些在多数情况下由品行极坏的主教付洗的人，你将说什么？这些人后来，当天主愿意时，因罪行被定罪，甚至被剥夺了职务本身，或完全被剥夺了共融。或者，对于那些可能由意见不正确或非常无知的主教付洗的人，你如何决定？当他们可能在传承圣事时没有清楚、诚实地说话，甚至说得不合宜，或者至少问了些问题，并从回答者那里听到了绝不应该如此询问或回答的事情时，情况又如何？尽管如此，这并不会大大损害我们真正的信德，尽管这些较为单纯的人可能没有用你那种优雅和秩序来传授信德的奥迹。你肯定会以你那惊人的谨慎说，这些人也应该重新受圣洗，因为这正是他们所欠缺的，或者阻碍了他们能够毫无玷污地接受那神圣而不可侵犯的信德奥迹。然而，卓越的人哪，让我们将能力归于天上的媒介，并将神圣威严的俯就允许其适当的运作；理解其中的巨大益处，让我们欣然接受它。如此，正如我们的救恩建立在圣神的圣洗之上（这圣神的圣洗多半与水的圣洗相关联），如果圣洗确实由我们施行，就应当完整而庄严地，并用所有那些记载的方式授予，且不缺少任何部分。或者，如果因情况需要，由较低级的圣职人员施行，就让我们等候结果，由我们补充，或由主保留来补充。然而，如果是由外人施行的，就应当尽可能并以允许的方式改正此事。因为在教会之外没有圣神，健全的信德也不能存在，不仅是在异端者中，甚至在那些处于分裂中的人中也是如此。因此，那些通过真理的教导以及他们自己的信德（这信德后来因心灵的净化而改善）而悔改并被纠正的人，只应借助于属神的圣洗，即通过主教的覆手和圣神的服务来帮助。此外，信德的完美印记在教会中按此方式并依此原则正确地授予。这样，对耶稣圣名的呼求（这是不可废除的）就不会被我们轻视；这当然是不合适的；尽管这样的呼求，如果它之后没有我们所提到的那些事情，可能会失效并失去救恩的效果。因为当宗徒说有\BibleQuote{一个洗礼}\BibleRef{弗 4:5}时，它必然因对耶稣圣名的持续呼求而发生，因为一旦呼求，任何人都不能取走它，即使我们敢于违背宗徒的决定，通过给予太多，甚至是通过增加圣洗的欲望来重复它。如果返回教会的人不愿再受圣洗，结果可能是我们剥夺了他圣神的圣洗，而我们却认为不应剥夺他水的圣洗。\par

11. 对于那个听了圣言，或许在基督的名下被接纳，立即宣认，并在被授予水的圣洗之前就受到惩罚的人，你将如何决定？你将宣称他因未受水的圣洗而丧亡了吗？或者，你确实会认为，即使他没有受水的圣洗，也有某种外在的东西帮助他得救？你认为他丧亡的想法将与主的话相悖，主说：\BibleQuote{凡在人前承认我的，我在我天上的父前也必承认他}；\BibleRef{玛 10:32} 因为无论为主而宣认者是听道者还是信徒，只要他宣认他应当宣认的那位基督，就无关紧要；因为主通过承认他，反过来以殉道的荣耀在他父前赐予恩宠给他的宣认者，正如他所应许的。但这当然不应被太宽泛地理解，仿佛可以扩展到如此程度，以至于任何异端者能宣认基督的名，尽管他实际上否认基督本身；他相信另一位基督，而基督宣称这对他毫无益处；因为主说，我们必须在我们面前宣认他，这没有他就不能做到，也没有对他名的崇敬。因此，两者都必须由宣认者持守：健全、真诚、无玷污、无侵害，而不对宣认者本身作任何选择，无论他是义人还是罪人，是完全的基督徒还是不完全的基督徒，只要他不怕在自己最大的危险中宣认主。这并不与前面的讨论相悖，因为其中留下了纠正许多恶事的时间，并且某些事只为了我们主的名而被允许；而殉道除非在主内并由主自己完成，否则不能圆满，因此没有人能在没有他的名的情况下宣认基督，基督的名也不能在没有基督自己的情况下为宣认者带来任何益处。\par

12. 因此，必须考虑这整个讨论，使之更加清楚。因为呼求耶稣之名只有随后得到适当补充才会有益处，因为先知和宗徒都如此宣告过。因为雅各伯在宗徒大事录中说：\Quote{诸位仁人弟兄，请听：西满述说了天主当初怎样眷顾外邦人，从他们中选取一个百姓归于自己的名下。先知的话也与此相符，正如经上所载：\InnerQuote{此后我要回来，重建达味已经倒塌的帐幕；我要重建其废墟，并要重新竖立它，使余剩的人寻求主，所有外邦人，凡我名被称于他们身上的，都寻求主——那行这些事的主说。}}\BibleRef{宗 15:13}因此，余剩的人——即部分犹太人以及所有主名被称于其上的外邦人——能够并且必须寻求主，因为正是那名字的呼求给他们提供了机会，甚至加给他们寻求主的必要性。并且与这些人一起，他们规定圣经——无论是全部还是仅一部分——可以用来比对外邦人更大胆地讨论关于真理的事；因为永生的天主之子、主耶稣的名字还没有被呼求在外邦人身上，同样也没有被呼求在只接受旧约经书的犹太人身上。因此，这两种人——即犹太人和外邦人——只要按当有的方式完全相信，就同样受洗。但是异端者虽然已经在水中以耶稣基督之名受过洗，却只应再以圣神受洗；并且在耶稣——即\Quote{天下间赐下唯一的名，使我们必须靠它得救}——中，死亡被合理地轻视，尽管如果他们保持现状，就不能得救，因为他们没有在他们身上被呼求主名之后寻求主——就像那些因着假基督而可能拒绝相信的人，关于他们主说：\Quote{你们要谨慎，免得有人迷惑你们。因为将有许多人假冒我的名来说：我是基督，并且要迷惑许多人。}祂又说：\Quote{那时若有人对你们说：看，基督在这里；或说：在那里；你们不要相信。因为将有假基督、假先知兴起，显大神迹和大奇事，以致如果可能，连被选的人也要被迷惑。}\BibleRef{玛 24:23}这些奇迹，毫无疑问，他们那时将在基督之名下行出来；在这个名下，甚至现在有些人似乎行了某些奇迹，并说假预言。但可以肯定的是，那些人，因为他们本身不属于基督，所以不属于基督；正如一个人若离开基督，只保留祂的名，他也不会得到多大益处；不，反而被那名所累，尽管他以前可能非常忠信，或非常正义，或受尊崇担任某个圣职，或享有宣认的尊荣。因为所有那些人，由于否认真正的基督，并引入或跟随另一位——虽然根本没有另一位——就不留给自己任何希望或救恩；不亚于那些在人前否认基督的人，他们必然被基督否认；对他们之前的品行、情感或尊荣不予考虑，正如他们本人敢于废除基督——即他们自己的救恩——一样，他们被这种简短的判决所定罪，因为主曾明确说过：\Quote{凡在人前否认我的，我在我天上的父前也必否认他。}正如\Quote{谁}这个词，也在宣认的语句中，最充分地给我们显示，宣认者自身的任何条件都不能成为障碍，尽管他以前可能是一个否认者，或异端者，或听者，或刚开始听的人，尚未受洗或未从异端皈依信仰真理，或是离开了教会后来又回来的人，然后当他回来时，还未等主教给他覆手，就被捕了，被迫在人前宣认基督；正如对于再次否认基督的人，任何特殊的昔日尊荣都不能为他生效以获得救恩。\par

13. 因为我们任何一个人都会认为，在这一点上，一个人在最后时刻所发现的一切，就是他必须受审判的依据，而他先前所做的一切都被抹去并消除了。因此，虽然在殉道中，在片刻之间就有如此大的变化，以致在最迅速的情况下，一切都可以改变；但任何人若失去了获得光荣救恩的机会，万一因自己的过错而将自己排除在外，就不要自欺欺人；正如\Person{罗特}的妻子，在类似的困境中，仅仅违背天使的命令回头一看，就变成了一根盐柱。同样，按照这个原则，那个异端者，若因宣认基督的名而被处死，如果他曾对天主或基督有任何错误的想法，尽管他因相信另一位天主或另一位基督而欺骗了自己，他后来也不能修正任何东西：他不是基督的宣认者，而只是名义上属于基督；因为宗徒接着说：\BibleQuote{如果我交出我的身体，使我被火焚烧，但没有爱德，我毫无益处。}因为藉此行为，那没有爱德的人——即没有爱那由法律、先知和福音如此宣告的天主和基督：\BibleQuote{你当全心、全意、全思爱主你的天主；你当爱近人如己。因为全部法律和先知都系于这两条诫命。}\BibleRef{玛 22:37}——同样，福音作者\Person{若望}说：\BibleQuote{凡爱的，都由天主而生，并且认识天主，因为天主是爱。}\BibleRef{若 4:7}同样，天主也说：\BibleQuote{天主如此爱了世界，以至于赐下了祂的独生子，使凡信祂的人不致丧亡，反而获得永生。}如此显然可见，那没有这爱德（即爱我们并被我们所爱）的人，从空洞的宣认和苦难中毫无收益，除非因此显明并证实他是个异端者，他相信另一位天主，或接受另一位基督，而非旧约和新约圣经所明确宣告的那位，这些圣经毫无隐晦地宣告全能的天父、万物的创造者，及祂的圣子。因为他们将如同那期待从另一位天主得救恩的人一样。最后，与他们所想相反，他们要被基督——全能天主圣父之子，他们所亵渎的创造主——定罪受永远的惩罚，那时天主要按照福音，藉基督耶稣审判人心中的隐密事，因为他们虽在祂的名内受洗，却不信祂。\par

14. 而且，甚至到了这个地步，整个那种异端者的圣洗也可以被修正，只要经过一段时间，如果这个人存活并修正了他的信仰，正如我们的天主在路加福音中对祂的门徒所说：\BibleQuote{但我还有当受的洗礼，}\BibleRef{路 12:50}同样，按照马尔谷福音的记载，祂怀着同一目的对\Person{载伯德}的儿子们说：\BibleQuote{你们能饮我所饮的杯吗？或者能受我所受的洗吗？}\BibleRef{谷 10:38}因为祂知道这些人不仅要受水的圣洗，也要受自己血的圣洗；这样，只受这一种圣洗的人，也能获得洗涤池的纯正信仰和纯朴的爱德，而受了两种圣洗的人，也同样地达到同一程度的救恩与光荣的圣洗。因为主所说的\BibleQuote{我还有当受的洗礼}在此处并非指第二个圣洗，仿佛有两个圣洗，而是表明还有另一种圣洗赐给了我们，与同一救恩相吻合。而且，这两种圣洗都应由我们的主亲自开创并祝圣，以便这两种中的一种或两种都能赐给我们这同一个具有双重的救恩与光荣的圣洗；并且同一圣洗的某些方式就这样为我们开启，以致有时其中一种可能缺失而无害，正如在聆听圣言的殉道者身上，水的圣洗缺失而无害；然而我们确信，这些人若有任何宽限，也习惯于受水的圣洗。同样，对于那些成为合法信徒的人，自己血的圣洗缺失也无害，因为他们在基督的名内受洗，已被主的最宝贵宝血所救赎；因为这两种主的圣洗之河都从同一源泉流出，使凡口渴的都可以来喝，正如圣经所说：\BibleQuote{从祂的腹中要流出活水的江河。}\BibleRef{若 7:38}这些河首先在主的受难中彰显出来，当时士兵用长枪刺透祂的肋旁，流出了血和水，以致同一人的同一肋旁涌出了两种不同的河流，使凡相信并饮这两条河的人，都充满圣神。因为主谈论这些河时，说明了这一点，暗示圣神就是那些信祂的人将要领受的：\BibleQuote{那时圣神还没有赐下，因为耶稣还没有受光荣。}\BibleRef{若 7:39}当祂如此说圣洗如何产生（宗徒宣布圣洗只有一个），那么根据这个原则显然可见，从同一伤口流出的水和血中产生了同一圣洗的不同种类；因为我们所谈到的两种水的圣洗，即同一类圣洗，虽然每种圣洗都应是单一的，正如我们已经更充分地论述的。\par

15. 既然我们似乎已将一切属神的圣洗以三重方式加以区分，现在让我们也来证明所提出的陈述，以免显得我们凭自己的判断贸然行事。因为若望在他的书信中教导我们，论及我们的主说：\Quote{这是那藉着水及血而来的耶稣基督；不单藉着水，而是藉着水及血；并且那作见证的是圣神，因为圣神是真理。原来作证的有三个：圣神、水及血，而这三个是一致的。}\BibleRef{若一 5:6}——好叫我们能从这些话中领悟：水常赐予\Emph{圣神}，人的血也常赐予圣神，圣神自己也常赐予圣神。因为水既如血般倾流，主也将圣神倾注于所有信者身上。诚然，人可以在水中受洗，同样也可在自己的血中受洗，而且尤其在圣神中受洗。因为伯多禄说：\Quote{这正是藉先知所预言的：在末日，天主说，我要将我的圣神倾注在所有血肉之人身上，你们的儿女要说预言，你们的青年要见异象，你们的老人要做异梦；在我的仆人和我的婢女身上，我也要倾注我的圣神。}\BibleRef{宗 2:17}——我们发现这圣神在旧约中已曾传达，虽非普遍或大量地，而是伴随着其他恩赐；或者，更是祂自愿降临于某些人，或覆盖他们，或临于他们，正如我们看见主对梅瑟论到七十位长老时所说的：\Quote{我要将降临在你身上的圣神，取一些分给他们。}\BibleRef{户 11:17}因此，按照祂的应许，天主将曾临于梅瑟的圣神的一部分放在他们身上，他们就在营中说出预言。梅瑟，作为一个属神的人，为这事的发生而欢喜，尽管他曾被纳维的儿子若苏厄不愿地劝说要反对这事，但并未因此被说服。此外，在民长纪和列王纪中，我们也看到圣神临于或降临于多人，如敖特尼耳、基德红、依弗大、三松、撒乌耳、达味，以及许多其他人。由此得出一个结论：主藉着他们最清楚地教导我们圣神依其自由意志接近的自由与能力，说：\Quote{圣神随意向哪里吹；你听到祂的响声，却不知道祂从哪里来，往哪里去。}因此，同样一位圣神有时也临于那些不配的人身上；这当然不是徒然或无理由的，而是为了一些必要的运作；如同圣神临于撒乌耳身上，他就说出预言。然而，在后来的日子里，主的圣神离开了他，并由一个来自主的恶神折磨他，因为他那时已带着之前小心差遣的使者前来，意图杀害达味；那些使者于是落入先知的行列中，也说出预言，以致他们既不能也不愿执行所受的命令。我们相信，临于他们众人身上的圣神，以天主所愿的奇妙智慧成就了这事。这圣神也充满了洗者若翰，甚至从他母胎起；又降临于百夫长科尔乃略同来的人身上，在他们尚未用水受洗之前。因此，圣神或是在人的圣洗之前，或是在其后，紧随着；或是在没有水的圣洗时，降临于信者身上。我们被劝告：要么我们应当恰当地保持圣洗的完整性，要么若偶然有人因耶稣基督之名施洗，我们应当予以补充，同时守护对耶稣基督圣名的最神圣的呼求，正如我们已极充分地陈明；并且守护那如此长久且由如此伟大人物所传承的习俗与权威。\par

16. 既然这论证的第一部分似乎已展开，我们应当触及随后的部分，这是为了异端者的缘故；因为很有必要不放过那偶尔落入我们手中的讨论，以免某个异端者以他的狡诈胆敢攻击我们那些较为单纯的弟兄。因为洗者若翰曾说我们必须在圣神和火中受洗，由他接着说了\Emph{和火}，有些绝望的人竟敢将他们的邪恶推行到如此地步，因此十分狡猾的人寻求如何借此败坏、玷污甚至废置圣洁的圣洗。这些人从西满魔术师那里得着他们观念的起源，以多种方式通过种种错误施行那观念；关于此人，西满伯多禄在宗徒大事录中说：\Quote{愿你的银子与你一同丧亡，因为你以为天主的恩宠可以用钱买得；在这项工作上，你没有分也没有权利，因为你的心在天主面前不正。}\BibleRef{宗 8:20}这样的人做这一切事，是出于欺骗那些较为单纯或好奇之人的欲望。他们中有些人试图争辩说，他们只施行一种健全而完美的圣洗，不像我们那样施行残缺不全的圣洗；他们声称这种圣洗是这样施行的：当他们一进入水中，火立刻出现在水面上。如果这能藉某种诡计实现——正如许多这类的诡计被归给阿那克西劳斯——或者如果这能藉某种自然方法发生，或者如果他们以为看见了这事，或者如果某个恶毒存在的作为和魔法毒素能从水中逼出火来——他们仍宣称如此欺骗和诡计是完美的圣洗；如果信者们被迫领受了这种圣洗，那么无疑他们必已失去了他们原有的圣洗。正如一个士兵在起了誓之后逃离自己的营帐，却在敌人完全不同的营帐里愿起另一类截然不同的誓言，显然他这样便从旧誓言中被解除了。\par

17. 此外，如果这样的人再次回到你们当中，你们必定会犹豫他是否能领受圣洗；然而，只要他悔改，你们仍当尽己所能地帮助这人。因为这种奸淫的，甚至是谋杀的洗礼，若有任何其他作者，那么它必定是由这些同一异端者为此同一谬误所编撰的一本书，题为\OriginalTerm{保禄的宣讲}{The Preaching of Paul}；在这本书中，违反正典的一切，你将发现基督承认自己的罪——尽管祂独自一人从未犯过罪——并且几乎被祂的母亲玛利亚勉强地接受了若翰的洗礼。还有，当祂受洗时，水上见到火，这在两部福音中均无记载。而且，过了这么长时间以后，伯多禄和保禄在耶路撒冷商议福音、互相考虑、争辩并最终安排应行之事以后，才彼此认识，仿佛是第一次见面；还有其他一些可耻荒谬的同类捏造之事；所有这些，你都会发现被收集在那本书里。但那些并非对圣神本性无知的人，理解有关火的言论是指圣神本身。因为在宗徒大事录中，按照我们主的同一许诺，在五旬节那天，当圣神降临在门徒身上，为使他们得以在祂内受洗时，有好像火的舌头显现，分坐在每人头上，为表明他们是以圣神和火受洗——也就是说，以那圣神受洗，祂是火，或如同火，正如那在荆棘中燃烧却没有烧毁荆棘的火；又如那作为天使之神的火，正如圣经所说：\BibleQuote{祂以风为使者，以火焰为仆役。}你们若效法祂，或与祂为伴、分享祂，便不必害怕任何火，甚至那在审判之日走在主前面、将要烧尽全世界的火也不例外，唯有那些以圣神和火受洗的人得免。\par

18. 圣神直至今日仍然是人眼看不见的，正如主所说：\BibleQuote{圣神随意吹拂；你不知道祂从哪里来，往哪里去。}\BibleRef{若 3:8} 然而，在信德的奥迹和属神之洗的开端，同一圣神曾明显显现，像火一样坐在门徒身上。此外，天开了，降在主身上如同鸽子；因为许多事，甚至几乎所有未来之事，都是显而易见的——尽管它们本来只是不可见的——如今也部分地、或有时、或以象征，显示给人的肉眼和不信，为坚强并确证我们的信德。但我也不应忽略福音所清楚宣布的事。因为我们主对瘫子说：\BibleQuote{放心吧，孩子，你的罪赦了，}\BibleRef{玛 9:2} 为表明人心是因信德而洁净，以获得随后而来的罪赦。那城里的罪妇也得了这样的罪赦，主对她说：\BibleQuote{你的罪赦了。} 当四周坐席的人开始私下议论：\BibleQuote{这赦罪的是谁？}\BibleRef{路 7:50}——因为经师和法利塞人曾针对瘫子恼怒地嘀咕——主对那妇女说：\BibleQuote{你的信德救了你，平安地去吧。}\BibleRef{路 7:50} 由这一切可以看出，心因信德而洁净，灵魂因圣神而被洗洁；此外，身体用水洗濯，而且藉着血我们更可立刻获得救恩的赏报。\par

19. 我认为我们已充分追述了洗者若翰的宣告，我们从那里开始我们的论述，当他向犹太人所说的话：\BibleQuote{我固然以水洗你们，为使你们悔改；但在我以后来的那一位，比我更强，我连解祂的鞋带也不配：祂要以圣神和火洗你们。}\BibleRef{路 3:16} 此外，我认为我们也妥善地整理了若望宗徒的教导，他说：\BibleQuote{作见证的有三样：圣神、水及血，而这三样是一致的。} 并且，若我没有记错，我们也解释了主所说的话：\BibleQuote{若翰固然以水施洗，但你们将要以圣神受洗。}\BibleRef{宗 1:5} 此外，我认为我们作为这习惯的缘由，已给出了并非微不足道的理由。我们要小心，尽管我们会在后文再论及此事，但不要让任何人以为我们只是在单一论点上挑起这场辩论；尽管这习惯本身，在敬畏天主而谦卑的人当中，也应当占有一席之地。\par

\SourceRecord{Translated by Robert Ernest Wallis.；Ante-Nicene Fathers；Vol. 5.；Edited by Alexander Roberts, James Donaldson, and A. Cleveland Coxe.；Buffalo, NY: Christian Literature Publishing Co.,；1886.；<http://www.newadvent.org/fathers/0515.htm>.}
